% Transcription diplomatique française de EGA IV, quatrième partie, pp. 345--361.
\clearpage
\oldpage[IV]{345}
\phantomsection
\addcontentsline{toc}{section}{Errata et Addenda (Liste 3)}
\section*{ERRATA ET ADDENDA}
\markboth{ERRATA ET ADDENDA, LISTE 3}{ERRATA ET ADDENDA, LISTE 3}
\begin{center}\large (Liste 3)\end{center}

\subsection*{A) Erreurs typographiques}

\paragraph{(II, 3.3.2).}
Ligne 1 de la p.~58, remplacer $\sh{S}$ par $\sh{M}$. Ligne 6 de la
p.~58, remplacer $S(n)$ par $M(n)$.

\paragraph{(II, 4.1.1).}
Ligne 8 de la p.~71, remplacer $\sh{O}$-$P$-Module par
$\sh{O}_P$-Module.

\paragraph{(II, 5.4.5).}
Ligne 6 de la p.~102, remplacer $u$ par $f\circ u$.

\paragraph{(II, 5.4.8).}
Avant la ligne 4 du bas de la p.~102, remplacer $f\times1_{S'}$ par
$f_{(Y')}$ dans le diagramme. Ligne 4 du bas de la p.~102, remplacer
$f\times1_{S'}$ par $f_{(Y')}$.

\paragraph{(II, 6.4.4).}
Ligne 1 du haut de la p.~122, remplacer normal par intégralement clos.

\paragraph{(II, 6.4.7).}
Ligne 14 du bas de la p.~122, remplacer dimension par rang.

\paragraph{(II, 6.5.2).}
Ligne 5 de la p.~127, remplacer $(\sh{B}|U_\lambda)$ par
$(B|U_\lambda)$. Ligne 15 de la p.~127, supprimer le premier
$(N_{\sh{B}/\sh{A}}\omega_{\lambda\mu})$.

\paragraph{($0_{\mathrm{III}}$, 8.1.11).}
Ligne 8 du bas de la p.~8, remplacer $h_X(u)$ par $h'_X(u)$.

\paragraph{($0_{\mathrm{III}}$, 9.1.2).}
Ligne 10 du bas de la p.~12, remplacer $\sh{F}$ par $X$.

\paragraph{($0_{\mathrm{III}}$, 10.2.1).}
Ligne 6 du bas de la p.~18, remplacer : « en outre si » par : « si en
outre ».

\paragraph{(III, 4.4.12).}
Ligne 10 de la p.~138, remplacer : chap.~V par : chap.~IV (8.11.2).

\paragraph{(III, 5.2.5).}
Ligne 12 de la p.~153, remplacer (I, 5.5.4) par (II, 5.5.4).

\paragraph{(III, 7.4.7).}
Ligne 15 de la p.~57, remplacer : tout $A$-module par : tout
$A$-module $M$.

\paragraph{(III, 7.8.5).}
Ligne 18 du bas de la p.~74, remplacer $F$ par $\sh{F}$.

\paragraph{(III, 7.8.9).}
Ligne 9 du bas de la p.~75, remplacer $M$ par $\sh{M}$.

\paragraph{(III, 7.9.11).}
Ligne 17 du bas de la p.~79, remplacer chap.~V par chap.~VI.

\paragraph{($0_{\mathrm{IV}}$, 16.2.2.2).}
Ligne 13 de la p.~26, remplacer $M'\cap\mathfrak q^nM'$ par
$M'\cap\mathfrak q^nM$.

\paragraph{($0_{\mathrm{IV}}$, 17.3.7).}
Ligne 16 du bas de la p.~50, remplacer (17.1.10) par (17.3.3).

\paragraph{($0_{\mathrm{IV}}$, 18.1.3).}
Dans le diagramme (18.1.3.1), $p_2$ doit être à côté de la flèche
$G\to F$.

\paragraph{($0_{\mathrm{IV}}$, 19.8.7).}
Lignes 7, 8 et 9 du bas de la p.~111, remplacer (21.5.5) par (21.5.3).

\paragraph{($0_{\mathrm{IV}}$, 19.8.8).}
Ligne 2 du bas de la p.~112, remplacer $A$-module par $B$-module.

\paragraph{($0_{\mathrm{IV}}$, 20.5.4).}
Ligne 5 du bas de la p.~130, remplacer $\Omega_{B'/A'}$ par
$\Omega^1_{B'/A'}$.

\paragraph{($0_{\mathrm{IV}}$, 20.7.8).}
Ligne 9 du bas de la p.~149, remplacer (18.4.3) par (18.5.4). Ligne 4
du bas, remplacer (18.4.1) par (18.5.1).

\paragraph{($0_{\mathrm{IV}}$, 21.5.3).}
Ligne 1 du bas de la p.~166, remplacer ramifiée par ramifié.

\paragraph{($0_{\mathrm{IV}}$, 21.7.1).}
Ligne 4 de la p.~170, remplacer (20.6.15.1) par (20.6.17.1).

\paragraph{($0_{\mathrm{IV}}$, 22.1.1).}
Ligne 12 de la p.~183, remplacer (20.6.20) par (20.6.10).

\oldpage[IV]{346}
\paragraph{($0_{\mathrm{IV}}$, 22.3.3).}
Ligne 13 du bas de la p.~192, remplacer $\widehat B_{\mathfrak p}$ par
$\widehat B_{\mathfrak q}$.

\paragraph{($0_{\mathrm{IV}}$, 22.7.3).}
Ligne 21 du bas de la p.~212, remplacer $\operatorname{Ker}(i)$ par
$\operatorname{Ker}(i_0)$. Ligne 12 du bas de la p.~212, remplacer
$(\mathfrak p/\mathfrak p^2)\otimes_L L$ par
$(\mathfrak p/\mathfrak p^2)\otimes_A L$.

\paragraph{($0_{\mathrm{IV}}$, 23.1.3).}
Ligne 4 de la p.~215, remplacer $x'^q\in K'$ par $x'\in K'$. Ligne 9
de la p.~215, remplacer $x'^q\in A'$ par $x'^q\in A$.

\paragraph{({\texorpdfstring{\protect\hyperlink{ega4.c04.statement.2-1-1}{IV, 2.1.1}}{IV, 2.1.1}}).}
Ligne 4 du bas de la p.~5, remplacer (III, 1.4.15.1) par
(III, 1.4.15.5).

\paragraph{({\texorpdfstring{\protect\hyperlink{ega4.c04.prop.2-4-3}{IV, 2.4.3}}{IV, 2.4.3}}).}
Ligne 12 de la p.~20, remplacer
$(X_{\mathrm{red}}\times_{Y_{\mathrm{red}}}Y_{\mathrm{red}})_{\mathrm{red}}$
par
$(X_{\mathrm{red}}\times_{Y_{\mathrm{red}}}Y'_{\mathrm{red}})_{\mathrm{red}}$.

\paragraph{({\texorpdfstring{\protect\hyperlink{ega4.c04.prop.2-6-1}{IV, 2.6.1}}{IV, 2.6.1}}).}
Ligne 10 de la p.~27, remplacer : $X'\to X_{(S')}$ par
$X'=X_{(S')}$. 

\paragraph{({\texorpdfstring{\protect\hyperlink{ega4.c04.prop.4-2-1}{IV, 4.2.1}}{IV, 4.2.1}}).}
Ligne 20 du bas de la p.~55, remplacer
$\operatorname{deg.tr.}_kE$ par $\operatorname{deg.tr.}_LE$ et
$\operatorname{deg.tr.}_kL(t)$ par $\operatorname{deg.tr.}_LL(t)$;
ligne 19 du bas, échanger « première » et « seconde ».

\paragraph{({\texorpdfstring{\protect\hyperlink{ega4.c04.lem.4-4-2}{IV, 4.4.2}}{IV, 4.4.2}}).}
Ligne 1 du bas de la p.~59, remplacer $X$ par $X'$.

\paragraph{({\texorpdfstring{\protect\hyperlink{ega4.c04.prop.4-7-3}{IV, 4.7.3}}{IV, 4.7.3}}).}
Ligne 10 du bas de la p.~75, remplacer préschéma par $k$-préschéma.

\paragraph{({\texorpdfstring{\protect\hyperlink{ega4.c04.prop.5-6-1}{IV, 5.6.1}}{IV, 5.6.1}}).}
Ligne 16 de la p.~98, remplacer (en indice) $k(\mathfrak p($ par
$k(\mathfrak p)$.

\paragraph{({\texorpdfstring{\protect\hyperlink{ega4.c04.prop.5-6-5}{IV, 5.6.5}}{IV, 5.6.5}}).}
Ligne 20 de la p.~99, supprimer la seconde parenthèse dans $k(\xi))$.

\paragraph{({\texorpdfstring{\protect\hyperlink{ega4.c04.lem.5-6-11-1}{IV, 5.6.11.1}}{IV, 5.6.11.1}}).}
Ligne 16 de la p.~102, remplacer $K$ par $\mathfrak K$.

\paragraph{({\texorpdfstring{\protect\hyperlink{ega4.c04.cor.5-11-2}{IV, 5.11.2}}{IV, 5.11.2}}).}
Note de bas de page de la p.~123, remplacer 5.11.2) par (5.11.2).

\paragraph{({\texorpdfstring{\protect\hyperlink{ega4.c04.prop.6-1-5}{IV, 6.1.5}}{IV, 6.1.5}}).}
Ligne 2 du bas de la p.~136, remplacer (5.5.1.2) par
($0$, 16.3.9.2).

\paragraph{({\texorpdfstring{\protect\hyperlink{ega4.c04.prop.6-3-1}{IV, 6.3.1}}{IV, 6.3.1}}).}
Ligne 9 de la p.~139, remplacer $y\in\mathfrak m$ par
$y\in\mathfrak n$.

\paragraph{({\texorpdfstring{\protect\hyperlink{ega4.c04.cor.6-7-2}{IV, 6.7.2}}{IV, 6.7.2}}).}
Ligne 18 de la p.~146, remplacer $\sh{O}'_{x'}$ par $\sh{O}_{x'}$.

\paragraph{({\texorpdfstring{\protect\hyperlink{ega4.c04.statement.6-10-1}{IV, 6.10.1}}{IV, 6.10.1}}).}
Ligne 13 de la p.~155, remplacer
$\sh{F}\otimes_{\sh{O}_Y}\sh{O}_Y$ par
$\sh{F}\otimes_{\sh{O}_X}\sh{O}_Y$.

\paragraph{({\texorpdfstring{\protect\hyperlink{ega4.c04.prop.6-14-1}{IV, 6.14.1}}{IV, 6.14.1}}).}
Ligne 21 de la p.~170, remplacer : n\textsuperscript{o}~3, prop.~12
par : n\textsuperscript{o}~2, prop.~8.

\paragraph{({\texorpdfstring{\protect\hyperlink{ega4.c04.statement.6-15-2}{IV, 6.15.2}}{IV, 6.15.2}}).}
Ligne 13 de la p.~177, remplacer $t(t-v)$ par $t^2(t-v)^2$.

\paragraph{({\texorpdfstring{\protect\hyperlink{ega4.c04.prop.6-15-7}{IV, 6.15.7}}{IV, 6.15.7}}).}
Ligne 6 du bas de la p.~179, remplacer $X\otimes_Ak'$ par
$X\otimes_kk'$.

\paragraph{({\texorpdfstring{\protect\hyperlink{ega4.c04.item.7-5-4}{IV, 7.5.4}}{IV, 7.5.4}}).}
Ligne 14 du bas de la p.~206, remplacer
$\operatorname{Spec}(\widehat B\otimes k)$ par
$\operatorname{Spec}(\widehat B\otimes_{\widehat A}k)$.

\paragraph{({\texorpdfstring{\protect\hyperlink{ega4.c04.item.7-8-4}{IV, 7.8.4}}{IV, 7.8.4}}).}
Ligne 10 du bas de la p.~216, remplacer « vérifiant »
({\hyperlink{ega4.c04.def.7-8-2}{7.8.2}}, (ii) et (iii)) par
« vérifiant ({\hyperlink{ega4.c04.def.7-8-2}{7.8.2}}, (ii) et (iii)) ».

\paragraph{(IV, 2\textsuperscript{e} Partie, Bibliographie).}
Ligne 6 de la p.~224, remplacer : \emph{Amer. J. of Math.}, par
\emph{Annals of Math.}

\paragraph{({\texorpdfstring{\protect\hyperlink{ega4.c04.cor.8-5-8}{IV, 8.5.8}}{IV, 8.5.8}}).}
Ligne 6 de la p.~24, remplacer $\operatorname{Ker}g_\mu$ à la fin de
la ligne par $\operatorname{Ker}g_\lambda$.

\paragraph{({\texorpdfstring{\protect\hyperlink{ega4.c04.prop.8-10-1}{IV, 8.10.1}}{IV, 8.10.1}}).}
Ligne 13 du bas de la p.~36, remplacer (2.3.11) par
{\hyperlink{ega4.c04.cor.8-2-11}{(8.2.11)}}.

\paragraph{({\texorpdfstring{\protect\hyperlink{ega4.c04.prop.10-2-1}{IV, 10.2.1}}{IV, 10.2.1}}).}
Ligne 14 du bas de la p.~98, remplacer $y$ par $Y$.

\paragraph{({\texorpdfstring{\protect\hyperlink{ega4.c04.statement.10-4-11}{IV, 10.4.11}}{IV, 10.4.11}}).}
Ligne 10 du bas de la p.~104, remplacer (17.9.7) par
{\hyperlink{ega4.c04.prop.17-9-6}{(17.9.6)}}.

\paragraph{({\texorpdfstring{\protect\hyperlink{ega4.c04.thm.11-2-9}{IV, 11.2.9}}{IV, 11.2.9}}).}
Ligne 3 du bas de la p.~126, remplacer $A_\lambda$ par $A_\mu$.

\paragraph{({\texorpdfstring{\protect\hyperlink{ega4.c04.prop.11-3-4}{IV, 11.3.4}}{IV, 11.3.4}}).}
Ligne 16 de la p.~134, remplacer $F$ par $\sh{F}$.

\paragraph{({\texorpdfstring{\protect\hyperlink{ega4.c04.thm.11-3-10}{IV, 11.3.10}}{IV, 11.3.10}}).}
Ligne 9 du bas de la p.~139, remplacer (11.3.7) par
{\hyperlink{ega4.c04.thm.11-3-10}{(11.3.10)}}.

\paragraph{({\texorpdfstring{\protect\hyperlink{ega4.c04.statement.11-7-5}{IV, 11.7.5}}{IV, 11.7.5}}).}
Ligne 11 de la p.~159, remplacer (6.5.11) par (6.15.11).

\paragraph{({\texorpdfstring{\protect\hyperlink{ega4.c04.thm.12-1-1}{IV, 12.1.1}}{IV, 12.1.1}}).}
Ligne 13 de la p.~174, ajouter à la fin : pour $k\geq1$.

\paragraph{({\texorpdfstring{\protect\hyperlink{ega4.c04.thm.14-2-1}{IV, 14.2.1}}{IV, 14.2.1}}).}
Ligne 11 de la p.~202, remplacer (6.1.2) par (6.1.1).

\oldpage[IV]{347}
\paragraph{({\texorpdfstring{\protect\hyperlink{ega4.c04.prop.14-3-8}{IV, 14.3.8}}{IV, 14.3.8}}).}
Ligne 10 du bas de la p.~206, remplacer $(f^{-1}y)$ par $f^{-1}(y)$.

\paragraph{({\texorpdfstring{\protect\hyperlink{ega4.c04.rem.14-5-7}{IV, 14.5.7}}{IV, 14.5.7}}).}
Ligne 16 de la p.~218, remplacer (14.5.5) par (14.5.6).

\paragraph{({\texorpdfstring{\protect\hyperlink{ega4.c04.cor.14-5-11}{IV, 14.5.11}}{IV, 14.5.11}}).}
Dans le diagramme du haut de la page 223, les flèches horizontales
doivent toutes aller de droite à gauche.

\paragraph{({\texorpdfstring{\protect\hyperlink{ega4.c04.prop.15-4-2}{IV, 15.4.2}}{IV, 15.4.2}}).}
Ligne 23 de la p.~230, remplacer (6.1.4) par (6.1.5).

\paragraph{({\texorpdfstring{\protect\hyperlink{ega4.c04.cor.15-5-2}{IV, 15.5.2}}{IV, 15.5.2}}).}
Ligne 16 de la p.~233, remplacer $f^{-1}(y)$ par $f^{-1}(y')$.

\paragraph{({\texorpdfstring{\protect\hyperlink{ega4.c04.item.15-6-10}{IV, 15.6.10}}{IV, 15.6.10}}).}
Ligne 6 de la p.~242, remplacer ${}^{-1}(y')$ par $f^{-1}(y')$.
Ligne 11 de la p.~242, remplacer : $\alpha$) et $\beta$), par :
$\alpha$) ou $\beta$).

\subsection*{B) Modifications de texte}

\paragraph{(Err\textsubscript{IV}, 1).}
Dans ($0_{\mathrm{III}}$, 11.1.6), l'énoncé (reproduisant un énoncé de
(G, I, 4.4.1)) est incorrect lorsque $q(n)$ est un entier quelconque;
il n'est correct que si l'on a $q(n)=0$ ou $q(n)=n$. En effet, dans le
premier cas, le noyau de $d_r^{p,0}$ est alors $E_r^{p,0}$, puisque
l'image de $d_r^{p,0}$ est $E_r^{p+r,-r+1}=0$ par hypothèse; de même
l'image de $d_r^{p-r,r-1}$ est nulle puisque
$E_r^{p-r,r-1}=0$; on a donc bien $E_{r+1}^{p,0}=E_r^{p,0}$. On
raisonne de même lorsque $q(n)=n$ pour tout $n$.

\paragraph{(Err\textsubscript{IV}, 2).}
Dans (III, 3.4), supprimer les lignes 2 et 3 de la p.~119.

\oldpage[IV]{348}
\paragraph{(Err\textsubscript{IV}, 3).}
Dans (III, 4.4.9), ligne 15 du bas, p.~137, remplacer : finies par :
finies en tant qu'ensembles.

\paragraph{(Err\textsubscript{IV}, 4).}
Dans (III, 7.8.10), ligne 9 de la p.~76, remplacer simple par lisse.

\paragraph{(Err\textsubscript{IV}, 5).}
Dans ($0_{\mathrm{IV}}$, 16.1.4), ligne 8 du bas de la p.~23, remplacer :
Dire que $A$ est caténaire, par : Dire qu'un anneau intègre $A$ est
caténaire.

\paragraph{(Err\textsubscript{IV}, 6).}
Dans ($0_{\mathrm{IV}}$, 16.1.9), on peut supprimer l'hypothèse que
$A$ est noethérien. En effet, pour tout élément $u\in B_M$, $A[u]$
est une sous-$A$-algèbre de $C$, donc $u$ est entier sur $A$ (Bourbaki,
\emph{Alg. comm.}, chap.~V, \S1, n\textsuperscript{o}~1, cor. de la
prop.~1), et par suite aussi sur $A_M$; on conclut encore par
{\hyperlink{ega4.c04.prop.16-1-5}{(16.1.5)}}.

\paragraph{(Err\textsubscript{IV}, 7).}
Dans ($0_{\mathrm{IV}}$, 16.5.12), on notera que la relation
(16.5.11.1) équivaut à la même relation où l'on remplace $A$ par les
anneaux quotients $A/\mathfrak r$, $\mathfrak r$ parcourant l'ensemble
des idéaux premiers minimaux de $A$. On peut donc dans la preuve de
{\hyperlink{ega4.c04.statement.16-5-12}{(16.5.12)}} se borner au cas
où $A$ est intègre, et appliquer (16.1.4.2).

\paragraph{(Err\textsubscript{IV}, 8).}
Supprimer ($0_{\mathrm{IV}}$, 19.3.12) dont la démonstration est erronée
(les ensembles $W_\lambda$, $U_\lambda$ et $V_\lambda$ ne sont pas des
variétés linéaires).

\paragraph{(Err\textsubscript{IV}, 9).}
Dans ($0_{\mathrm{IV}}$, 19.8.3), lignes 6 et 8 du bas de la p.~110,
remplacer homomorphisme par homomorphisme local; de même dans
($0_{\mathrm{IV}}$, 19.8.4), ligne 1 de la p.~111.

\paragraph{(Err\textsubscript{IV}, 10).}
Dans ($0_{\mathrm{IV}}$, 20.1.2), ligne 2 du bas de la p.~117,
remplacer $a\in A$ par $b\in B$ et $aD$ par $bD$; ligne 1 du bas de la
p.~117, remplacer $A$-module par $B$-module.

\paragraph{(Err\textsubscript{IV}, 11).}
Dans ($0_{\mathrm{IV}}$, 20.2.1), ligne 8 de la p.~119, remplacer
$A$-modules par $B$-modules. Lignes 21 et 22 de la p.~119, remplacer
« di-homomorphismes de modules (relatif à $u$) » par : homomorphisme
de $B$-modules.

\paragraph{(Err\textsubscript{IV}, 12).}
Dans ($0_{\mathrm{IV}}$, 20.4.6.1), remplacer $p_2-p_1$ par $p_1-p_2$;
ligne 10 du bas de la p.~125, remplacer $x\otimes1-1\otimes x$ par
$1\otimes x-x\otimes1$.

\paragraph{(Err\textsubscript{IV}, 13).}
Dans ($0_{\mathrm{IV}}$, 20.4.8), ligne 14 de la p.~126, remplacer
$p_2=p_1-d_{B/A}$ par $p_1=p_2-d_{B/A}$.

\paragraph{(Err\textsubscript{IV}, 14).}
Dans ($0_{\mathrm{IV}}$, 20.4.13), remplacer les lignes 6 à 11 du bas
de la p.~127 par le texte suivant :

\medskip
\noindent
(iv) Soient $A$, $B$ deux anneaux intègres tels que $A\subset B$, $B$
étant entier sur $A$, et que le corps des fractions de $B$ soit une
extension séparable de celui de $A$. Alors le $B$-module
$\Omega^1_{B/A}$ est un module de torsion. En effet, pour tout $x\in B$,
soit $f$ le polynôme minimal de $x$ par rapport au corps des fractions
$K$ de $A$; il existe $a\in A$ non nul tel que $af[T]\in A[T]$; comme
$x$ est séparable sur $K$, on a $f'(x)\neq0$, et d'autre part de la
relation $af(x)=0$ on déduit $af'(x)d_{B/A}x=0$, d'où notre assertion
en vertu de {\hyperlink{ega4.c04.cor.20-4-7}{(20.4.7)}}.

\paragraph{(Err\textsubscript{IV}, 15).}
Dans ($0_{\mathrm{IV}}$, 22.2.7), ligne 5 de la p.~189, au lieu de
« formellement lisses », il faut « formellement lisses pour les
topologies préadiques ».

\paragraph{(Err\textsubscript{IV}, 16).}
Dans le Sommaire du chap.~IV, ligne 4 du bas de la p.~222, ajouter :
et anneaux strictement locaux. Lignes 1 à 3 du bas, remplacer par :
\begin{quote}
\S19. Immersions régulières et platitude normale.\\
\S20. Fonctions méromorphes et pseudo-morphismes.\\
\S21. Diviseurs.
\end{quote}
Ligne 16 de la p.~223, remplacer « de présentation finie » par
« localement de présentation finie »; ligne 19 de la p.~223, remplacer
« comme » par « comme étant ».

\paragraph{(Err\textsubscript{IV}, 17).}
Dans (IV, 3.1.10), ajouter après la ligne 16 de la p.~39 :

\begin{remark}[3.1.10.2]
\label{IV.3.1.10.2}
Soit $f:X\to Y$ un morphisme quasi-compact et quasi-séparé, et soit
$\sh{F}$ un $\sh{O}_X$-Module quasi-cohérent, de sorte que
$f_*(\sh{F})$ est un $\sh{O}_Y$-Module quasi-cohérent
{\hyperlink{ega4.c04.statement.1-7-4}{(1.7.4)}}. Alors on a
\begin{align}
 \operatorname{Ass}(f_*(\sh{F}))&\subset f(X),
 \label{IV.3.1.10.3}\tag{3.1.10.3}\\
 \operatorname{Ass}(f_*(\sh{F}))&\subset f(\operatorname{Ass}(\sh{F}))
 \quad\text{si $X$ est localement noethérien}.
 \label{IV.3.1.10.4}\tag{3.1.10.4}
\end{align}

Soit en effet $y\in\operatorname{Ass}(f_*(\sh{F}))$; en vertu de
{\hyperlink{ega4.c04.lem.2-3-1}{(2.3.1)}} et du fait que le morphisme
$\operatorname{Spec}(\sh{O}_{Y,y})\to Y$ est plat,
$(f_*(\sh{F}))_y$ est égal à $(f'_*(\sh{F}'))_y$, où
\[
 f':X\times_Y\operatorname{Spec}(\sh{O}_{Y,y})
 \longrightarrow\operatorname{Spec}(\sh{O}_{Y,y})
\]
est déduit de $f$ par changement de base et
$\sh{F}'=\sh{F}\otimes_{\sh{O}_Y}\sh{O}_{Y,y}$. On est donc ramené à
prouver (3.1.10.3) et (3.1.10.4) en y remplaçant $Y$ par
$\operatorname{Spec}(\sh{O}_{Y,y})$, $f$ par $f'$ et $\sh{F}$ par
$\sh{F}'$, compte tenu de (I, 3.6.5) et de la définition de
$\operatorname{Ass}(\sh{F}')$.

Par définition, il y a un élément
$t_y\in(f_*(\sh{F}))_y$ non nul et dont $\mathfrak m_y$ est
l'annulateur, et comme $(f_*(\sh{F}))_y$ s'identifie à
$\Gamma(Y,f_*(\sh{F}))$ puisque $Y$ est local, $t_y$ s'identifie aussi
à une section $s\in\Gamma(X,\sh{F})$. L'Idéal de $\sh{O}_X$ annulateur
de $s$ contient donc $\mathfrak m_y\sh{O}_X$, et il résulte de
($0_{\mathrm I}$, 1.7.4) que le support de $s$ est contenu dans
$f^{-1}(y)$; si ce dernier était vide, on aurait $s=0$, ce qui est
absurde. Lorsque de plus $X$ est localement noethérien, alors le
sous-$\sh{O}_X$-Module cohérent $\sh{O}_Xs$ de $\sh{F}$ n'est pas nul,
donc $\operatorname{Ass}(\sh{O}_Xs)\neq\varnothing$
{\hyperlink{ega4.c04.cor.3-1-5}{(3.1.5)}}; mais
\[
 \operatorname{Ass}(\sh{O}_Xs)\subset
 \operatorname{Supp}(\sh{O}_Xs)\subset f^{-1}(y),\qquad
 \operatorname{Ass}(\sh{O}_Xs)\subset\operatorname{Ass}(\sh{F})
\]
{\hyperlink{ega4.c04.prop.3-1-7}{(3.1.7)}}, donc
$y\in f(\operatorname{Ass}(\sh{F}))$.
\end{remark}

\oldpage[IV]{349}
\begin{remark}[3.1.10.5]
\label{IV.3.1.10.5}
On notera que même lorsque $f$ est un morphisme propre, on n'a pas
nécessairement
\[
 f(\operatorname{Ass}(\sh{F}))\subset
 \operatorname{Ass}(f_*(\sh{F})).
\]
On le voit en prenant pour $Y$ le spectre d'un corps et pour $X$ un
espace projectif sur $k$, par exemple $\mathbf P^1_k$; alors
$\sh{F}=\sh{O}_X(-1)$ est un $\sh{O}_X$-Module cohérent $\neq0$ tel
que $\Gamma(X,\sh{F})=0$ (III, 2.1.13), donc
$\operatorname{Ass}(f_*(\sh{F}))=\varnothing$ et
$\operatorname{Ass}(\sh{F})\neq\varnothing$.
\end{remark}

\paragraph{(Err\textsubscript{IV}, 18).}
Après (IV, 3.2.8), insérer après la ligne 17 de la p.~43 :

\begin{remark}[3.2.9]
\label{IV.3.2.9}
Pour tout point $s$ d'un préschéma $X$, notons
$i_s:\operatorname{Spec}(k(s))\to X$ le morphisme canonique, qui est
un monomorphisme quasi-compact. Pour tout $k(s)$-module $G(s)$,
considéré comme Module sur $\operatorname{Spec}(k(s))$,
$(i_s)_*(G(s))$ est donc un $\sh{O}_X$-Module quasi-cohérent
{\hyperlink{ega4.c04.statement.1-7-4}{(1.7.4)}}. On vérifie aussitôt
que la fibre de $(i_s)_*(G(s))$ est nulle en tout point
$x\notin\overline{\{s\}}$ et isomorphe à $G(s)$ dans le cas contraire.
Si $X$ est localement noethérien et $G(s)$ de longueur finie, il est
clair que $\operatorname{Ass}((i_s)_*(G(s)))$ est réduit au seul point
$s$, autrement dit $(i_s)_*(G(s))$ est irredondant.

Cela étant, supposons $X$ localement noethérien, et soit $\sh{F}$ un
$\sh{O}_X$-Module cohérent. Soit d'autre part $S$ une partie discrète
de $X$, et pour tout $s\in S$, posons
$\sh{F}(s)=i_s^*(\sh{F})$, qui est un $k(s)$-module; considérons pour
chaque $s\in S$ un module quotient $G(s)$ de $\sh{F}(s)$ de longueur
finie. On a donc un homomorphisme canonique
\[
 u:\sh{F}\longrightarrow
 \bigoplus_{s\in S}(i_s)_*(\sh{F}(s))\longrightarrow
 \bigoplus_{s\in S}(i_s)_*(G(s))
\]
de $\sh{O}_X$-Modules cohérents (puisque $S$, étant discret, est
localement fini dans $X$). Si $\sh{H}=\operatorname{Ker}u$,
$\sh{G}=\sh{F}/\sh{H}$ est un $\sh{O}_X$-Module cohérent et on déduit
de $u$ un homomorphisme injectif
\[
 v:\sh{G}\longrightarrow\bigoplus_{s\in S}(i_s)_*(G(s)).
\]
Si $\sh{G}^{(s)}$ est l'image de $\sh{G}$ dans $(i_s)_*(G(s))$,
$(\sh{G}^{(s)})_{s\in S}$ est une décomposition irredondante, qui est
réduite lorsque les $G(s)$ sont tous $\neq0$, et telle que
$\operatorname{Ass}(\sh{G})=S$ ($\sh{G}$ étant donc sans cycle premier
associé immergé). On a ainsi obtenu une bijection
$(G(s))_{s\in S}\leftrightarrow\sh{G}$, où $S$ parcourt l'ensemble des
parties discrètes de $X$ pour lesquelles $\sh{F}(s)\neq0$ pour tout
$s\in S$, et, pour chaque $S$, $(G(s))_{s\in S}$ parcourt l'ensemble
des quotients $\neq0$ et de longueur finie des $\sh{F}(s)$; d'autre
part $\sh{G}$ parcourt l'ensemble des $\sh{O}_X$-Modules quotients
cohérents de $\sh{F}$, sans cycle premier immergé. La bijection
réciproque associe à chaque $\sh{G}$ sa décomposition irredondante
réduite.
\end{remark}

\paragraph{(Err\textsubscript{IV}, 19).}
Dans (IV, 4.2.1), avant la ligne 17 du bas de la p.~55, insérer :

\begin{remark}[4.2.1.4]
\label{IV.4.2.1.4}
Sous les hypothèses de {\hyperlink{ega4.c04.prop.4-2-1}{(4.2.1)}}, on a
en outre
\begin{equation}
 \dim(K\otimes_kL)=
 \inf(\operatorname{deg.tr.}_kK,\operatorname{deg.tr.}_kL).
 \label{IV.4.2.1.5}\tag{4.2.1.5}
\end{equation}
où il doit être entendu que si les deux corps $K$, $L$ sont de degré de
transcendance infini (qui peuvent être des cardinaux quelconques), le
premier membre doit être pris égal à $+\infty$ (la dimension considérée
ici est celle définie dans ($0$, 16.1.1)).

Supposons en effet par exemple que
$\operatorname{deg.tr.}_kK\leq\operatorname{deg.tr.}_kL$ et en premier
lieu que $\operatorname{deg.tr.}_kK=n<+\infty$. On a vu dans la preuve
de {\hyperlink{ega4.c04.prop.4-2-1}{(4.2.1)}} que $K\otimes_kL$ est une
extension algébrique de $K'\otimes_kL$, et a donc même dimension
($0$, 16.1.5). On peut donc se borner au cas où
$K=k[T_1,\ldots,T_n]$; alors $K\otimes_kL$ est un anneau de fractions
de
\[
 k[T_1,\ldots,T_n]\otimes_kL=L[T_1,\ldots,T_n],
\]
qui est de dimension $n$ {\hyperlink{ega4.c04.cor.5-5-4}{(5.5.4)}}, et
par suite $\dim(K\otimes_kL)\leq n$ ($0$, 16.1.3.1). Pour prouver que
$\dim(K\otimes_kL)=n$, il suffit évidemment de montrer qu'il y a dans
$\operatorname{Spec}(L[T_1,\ldots,T_n])$ un point rationnel sur $L$
qui appartient à $\operatorname{Spec}(K\otimes_kL)$
{\hyperlink{ega4.c04.cor.5-2-3}{(5.2.3)}}, ou encore (I, 1.7.4) qu'il
existe un $L$-homomorphisme $u:K\otimes_kL\to L$. Mais par définition
du produit tensoriel, cela revient à prouver l'existence d'un
$k$-homomorphisme $v:K\to L$, ce qui est évident en vertu de
l'hypothèse
$\operatorname{deg.tr.}_kK\leq\operatorname{deg.tr.}_kL$.

Si $\operatorname{deg.tr.}_kK\leq\operatorname{deg.tr.}_kL$ et si ces
deux cardinaux sont infinis, on voit de la même façon qu'il y a un point
rationnel dans $\operatorname{Spec}(L[(t_\alpha)_{\alpha\in I}])$ qui
appartient à $\operatorname{Spec}(K\otimes_kL)$; mais le noyau de tout
$L$-homomorphisme de $L[(t_\alpha)_{\alpha\in I}]$ dans $L$ est un
idéal maximal engendré par une famille de la forme
$(t_\alpha-a_\alpha)_{\alpha\in I}$ avec $a_\alpha\in L$. Il est clair
qu'un tel idéal maximal contient une suite infinie décroissante
d'idéaux premiers distincts, donc l'anneau local de
$L[(t_\alpha)_{\alpha\in I}]$ correspondant est non noethérien et de
dimension infinie; comme c'est un anneau local de $K\otimes_kL$, on
voit que ce dernier est non noethérien et de dimension infinie.
\end{remark}

\oldpage[IV]{350}
\paragraph{(Err\textsubscript{IV}, 20).}
Après (IV, 4.5.19), insérer, avant la ligne 12 du bas de la p.~67 :

\begin{corollary}[4.5.19.3]
\label{IV.4.5.19.3}
Soient $k$ un corps, $X$ un $k$-préschéma de type fini sur $k$, qui est
irréductible mais non géométriquement irréductible. Alors il existe une
partie fermée $Y$ de $X$, rare dans $X$, contenant tous les points
$x\in X$ tels que $k(x)$ soit radiciel sur $k$ (et en particulier tous
les points de $X$ rationnels sur $k$).
\end{corollary}

Soient $k'$ une clôture algébrique de $k$, $X'=X\otimes_kk'$,
$p:X'\to X$ la projection canonique. Si $x\in X$ est tel que $k(x)$
soit radiciel sur $k$, il n'existe qu'un seul point $x'\in X'$ au-dessus
de $x$ (I, 3.4.9); ce point appartient nécessairement à deux composantes
irréductibles distinctes de $X'$, sans quoi $X$ serait géométriquement
irréductible en vertu de
{\hyperlink{ega4.c04.prop.4-5-19}{(4.5.19)}}, contrairement à
l'hypothèse. Soient alors $X'_i$ ($1\leq i\leq n$) les composantes
irréductibles du préschéma noethérien $X'$, et pour
$1\leq i<j\leq n$, posons $X'_{ij}=X'_i\cap X'_j$; l'ensemble
\[
 Y=p\left(\bigcup_{i<j}X'_{ij}\right)
\]
répond à la question. En effet, comme $p$ est un morphisme entier et
que les $X'_{ij}$ sont fermés dans $X'$, $Y$ est fermé dans $X$.
D'autre part, si $\xi$ est le point générique de $X$, $p^{-1}(\xi)$
est discret et formé des points génériques des $X'_i$
{\hyperlink{ega4.c04.cor.4-5-11}{(4.5.11)}}, donc ne rencontre aucun
des $X'_{ij}$; l'ensemble fermé $Y$ est par suite rare.

\paragraph{(Err\textsubscript{IV}, 21).}
Dans (IV, 5.6.11.1), ligne 7 du bas de la p.~102, après « intègre »,
ajouter : « de dimension 2 ».

\paragraph{(Err\textsubscript{IV}, 22).}
Après (IV, 6.2.3), après la ligne 18 de la p.~138, insérer :

\begin{proposition}[6.2.4]
\label{IV.6.2.4}
Sous les hypothèses de {\hyperlink{ega4.c04.prop.6-2-3}{(6.2.3)}}, soit
de plus $M$ un $A$-module de type fini et supposons $M\neq0$ et
$N\neq0$. Alors on a
\begin{equation}
 \operatorname{dim.proj}_B(M\otimes_AN)=
 \operatorname{dim.proj}_A(M)+
 \operatorname{dim.proj}_{B\otimes_Ak}(N\otimes_Ak).
 \label{IV.6.2.4.1}\tag{6.2.4.1}
\end{equation}
\end{proposition}

Considérons la résolution gauche $L_\bullet$ de $N$ introduite dans
{\hyperlink{ega4.c04.prop.6-2-3}{(6.2.3)}}, et une résolution gauche
$K_\bullet$ de $M$ par des $A$-modules libres de type fini :
\[
 \cdots\longrightarrow K_i\longrightarrow K_{i-1}\longrightarrow
 \cdots\longrightarrow K_0\longrightarrow M\longrightarrow0.
\]
On a vu que les $L_i$, $N$, les $Z_i(L_\bullet)$ sont des $A$-modules
plats; il en est de même des $B_i(L_\bullet)$ qui sont égaux aux
$Z_i(L_\bullet)$ sauf pour $i=0$, et dans ce dernier cas
$B_0(L_\bullet)$ est

\oldpage[IV]{351}
plat par ($0_{\mathrm I}$, 6.1.2); enfin, les $H_i(L_\bullet)$ sont nuls
par définition sauf $H_0(L_\bullet)=N$, qui est un $A$-module plat. On
sait que, dans ces conditions,
\[
 H_n(K_\bullet\otimes_AL_\bullet)\simeq
 \bigoplus_{p+q=n}
 H_p(K_\bullet)\otimes_AH_q(L_\bullet)
\]
(G, I, 5.5), donc $K_\bullet\otimes_AL_\bullet$ (gradué comme
d'ordinaire) est une résolution de $M\otimes_AN$ par des $B$-modules
libres. Posons
\[
 r=\operatorname{dim.proj}_A(M),\qquad
 s=\operatorname{dim.proj}_B(N)=
 \operatorname{dim.proj}_{B\otimes_Ak}(N\otimes_Ak).
\]
Il y a donc une résolution libre $K_\bullet$ de longueur $r$ et une
résolution libre $L_\bullet$ de longueur $s$, et
$K_\bullet\otimes_AL_\bullet$ est alors une résolution libre de longueur
$r+s$, ce qui prouve déjà que
$\operatorname{dim.proj}_B(M\otimes_AN)\leq r+s$. Il reste à prouver
($0$, 17.2.6) que l'on a
\[
 \operatorname{Tor}^{B}_{r+s}(M\otimes_AN,k(B))\neq0.
\]
Utilisant la résolution précédente $K_\bullet\otimes_AL_\bullet$ pour
calculer ce foncteur, on obtient
\[
 \operatorname{Tor}^{B}_{r+s}(M\otimes_AN,k(B))=
 H_{r+s}((K_\bullet\otimes_AL_\bullet)\otimes_Bk(B)).
\]
Mais on peut écrire
\[
 (K_\bullet\otimes_AL_\bullet)\otimes_Bk(B)=
 (K_\bullet\otimes_Ak)\otimes_k
 (L_\bullet\otimes_Bk(B)),
\]
qui est donc un produit tensoriel de deux complexes d'espaces
vectoriels sur $k$; tous les $k$-modules étant plats, la formule de
Künneth donne
\[
 H_{r+s}((K_\bullet\otimes_Ak)\otimes_k
 (L_\bullet\otimes_Bk(B)))\simeq
 H_r(K_\bullet\otimes_Ak)\otimes_k
 H_s(L_\bullet\otimes_Bk(B)),
\]
l'homologie du premier (resp. second) complexe étant nulle en degrés
$>r$ (resp. $>s$). On a donc
\[
 \operatorname{Tor}^{B}_{r+s}(M\otimes_AN,k(B))=
 \operatorname{Tor}^{A}_{r}(M,k)\otimes_k
 \operatorname{Tor}^{B}_{s}(N,k(B))
\]
à isomorphie près, et comme chacun des facteurs est $\neq0$ par
hypothèse, il en est de même de leur produit tensoriel sur $k$.

\paragraph{(Err\textsubscript{IV}, 23).}
Dans (IV, 6.10.1), ligne 8 de la p.~155, après : canonique, ajouter :
$\sh{F}$ un $\sh{O}_X$-Module cohérent.

\paragraph{(Err\textsubscript{IV}, 24).}
Dans ({\hyperlink{ega4.c04.prop.6-14-4}{IV, 6.14.4}}), la référence
donnée ligne 12 du bas de la p.~175 est insuffisante; il faut raisonner
de la façon suivante. Tout d'abord, si $M$ est une base de transcendance
de $K'$ sur $A$, $A(M)$ est algébriquement fermé dans $B(A(M))$
(Bourbaki, \emph{Alg.}, chap.~V, \S6, exerc.~8 c)). On peut donc se
borner au cas où $K'$ est une extension algébrique séparable de $A$.
Soit $p$ l'exposant caractéristique de $A$; il résulte de Bourbaki,
\emph{loc. cit.}, \S9, exerc.~2, que la fermeture algébrique de $K'$
dans $B(K')$ est radicielle sur $K'$; montrons que tout élément
$x\in B(K')$ tel que $x^{p^r}\in K'$ appartient déjà à $K'$. Pour cela,
considérons une base $(e_\lambda)$ de $K'$ sur $A$; puisque $K'$ est une
extension algébrique séparable de $A$, les $e_\lambda^{p^r}$ forment
aussi une base de $K'$ sur $A$ (Bourbaki, \emph{loc. cit.}, \S8,
n\textsuperscript{o}~3, cor. de la prop.~4). On a
$x=\sum_\lambda e_\lambda b_\lambda$, avec $b_\lambda\in B$ et, par
hypothèse, il existe des éléments $a_\lambda\in A$ tels que l'on ait
\[
 \sum_\lambda e_\lambda^{p^r}b_\lambda^{p^r}=
 \sum_\lambda e_\lambda^{p^r}a_\lambda;
\]
mais cela entraîne $b_\lambda^{p^r}=a_\lambda$ pour tout $\lambda$;
comme par hypothèse $A$ est algébriquement fermé dans $B$, on a donc
$b_\lambda\in A$ pour tout $\lambda$, et par suite $x\in K'$.

\paragraph{(Err\textsubscript{IV}, 25).}
Dans (IV, 7.4.8, B)), p.~203, la question a été résolue affirmativement
par R.~Kiehl, \emph{Ausgezeichnete Ringe der nichtarchimedischen
analytischen Geometrie}, à paraître dans \emph{Inventiones
Mathematicae}.

\paragraph{(Err\textsubscript{IV}, 26).}
Dans (IV, 7.8.4), ligne 9 du bas de la p.~217, remplacer :
\begin{quote}
Nous verrons plus loin (\S18) que si $k$ est un corps valué complet non
discret\ldots
\end{quote}
par :
\begin{quote}
Il a été démontré par R.~Kiehl (dans l'article cité en
(Err\textsubscript{IV}, 25)) que si $k$ est un corps valué complet non
discret\ldots
\end{quote}

\oldpage[IV]{352}
\paragraph{(Err\textsubscript{IV}, 27).}
Dans (IV, 7.8.4), ajouter, avant la ligne 6 du bas de la p.~217 :

\smallskip
\noindent
\textup{(vi)} Soit $A$ un anneau local noethérien intègre de dimension
1; il est alors universellement caténaire
{\hyperlink{ega4.c04.cor.7-2-9}{(7.2.9)}} et, pour qu'il soit excellent,
il faut et il suffit donc, en vertu de
{\hyperlink{ega4.c04.statement.7-3-19}{(7.3.19)}}, que ses fibres
formelles soient géométriquement réduites, ou encore que son complété
$\widehat A$ soit réduit et que les corps composants de son anneau total
de fractions soient des extensions séparables du corps des fractions
$K$ de $A$. En vertu de {\hyperlink{ega4.c04.thm.7-6-4}{(7.6.4)}} et
{\hyperlink{ega4.c04.thm.7-7-2}{(7.7.2)}}, il revient aussi au même de
dire que $A$ est universellement japonais. Lorsque $A$ est un anneau de
valuation discrète, il en est de même de $\widehat A$; dire que $A$ est
excellent signifie alors que le corps des fractions de $\widehat A$
est une extension séparable de $K$, condition toujours satisfaite si
$K$ est de caractéristique 0.

\paragraph{(Err\textsubscript{IV}, 28).}
Après (IV, 8.10.3), après la ligne 14 de la p.~37, insérer :

\begin{remark}[8.10.3.1]
\label{IV.8.10.3.1}
On ne peut dans {\hyperlink{ega4.c04.prop.8-10-3}{(8.10.3)}} supprimer
l'hypothèse que les $\overline{f_\alpha(X_\alpha)}$ soient
constructibles. Prenons par exemple $S=\operatorname{Spec}(A)$, où $A$
est un anneau tel que $S$ soit compact, totalement discontinu et sans
point isolé, tout idéal premier $\mathfrak p$ de $A$ étant maximal et
$A_\mathfrak p$ étant un corps (Bourbaki, \emph{Alg. comm.}, chap.~II,
\S4, exerc.~16 et 17). Soit $s\in S$, et prenons pour
$Y_\alpha=S_\alpha$ les voisinages ouverts affines de $s$ dans $S$; on
prend $X_\alpha=\operatorname{Spec}(k(s))$ pour tout $\alpha$,
$f_\alpha$ et $f$ étant les morphismes canoniques; il est clair que
$f$, qui est l'identité, est dominant, mais aucun des $f_\alpha$ ne
l'est.
\end{remark}

\paragraph{(Err\textsubscript{IV}, 29).}
Dans (IV, 8.12.8), ligne 7 de la p.~47, remplacer « tel qu'il existe un
$\sh{O}_Y$-Module ample » par : « et quasi-séparé ». Ligne 12 de la
p.~47, remplacer « L'hypothèse » par : « Il résulte de
{\hyperlink{ega4.c04.prop.8-12-3}{(8.12.3)}} que l'existence de la
factorisation de l'énoncé est une propriété locale sur $Y$; on peut
donc supposer $Y$ affine, ce qui ».

\paragraph{(Err\textsubscript{IV}, 30).}
Après (IV, 8.12.11), ajouter à la fin de la p.~48 :

\begin{corollary}[8.12.12]
\label{IV.8.12.12}
Soient $X$, $Y$ deux préschémas intègres,
$f:X\to Y$ un morphisme séparé, birationnel et localement de type fini.
Soit $y\in Y$ un point normal de $Y$ et supposons qu'un point
$x\in f^{-1}(y)$ soit isolé dans $f^{-1}(y)$. Alors il existe un
voisinage ouvert $U$ de $y$ dans $Y$ tel que la restriction
$f^{-1}(U)\to U$ de $f$ soit un isomorphisme.
\end{corollary}

\begin{proof}
Il suffit de prouver qu'il existe un voisinage ouvert $U$ de $y$ dans
$Y$ et une $U$-section $g:U\to f^{-1}(U)$ de $f^{-1}(U)$ telle que
$g(y)=x$. En effet, comme $U$ contient le point générique $\eta$ de
$Y$ et $f^{-1}(U)$ le point générique $\xi$ de $X$, on a
nécessairement $g(\eta)=\xi$ puisque $f$ est birationnel
{\hyperlink{ega4.c04.statement.6-15-4}{(6.15.4)}}; comme $g$ est une
immersion, $g(U)$ est localement fermé dans $f^{-1}(U)$, donc ouvert
dans son adhérence dans $f^{-1}(U)$, et comme cette dernière est égale
à $f^{-1}(U)$ d'après ce qui précède, $g(U)$ est ouvert dans
$f^{-1}(U)$; comme le sous-préschéma de $f^{-1}(U)$ associé à $g$ est
isomorphe à $U$, donc réduit, et que $f^{-1}(U)$ est réduit, il résulte
de (I, 5.2.1) que $g$ est une immersion ouverte. D'autre part, puisque
$f$ est séparé, $g$ est une immersion fermée (I, 5.4.6), donc $g(U)$
est un ensemble à la fois ouvert et fermé dans $f^{-1}(U)$. Mais comme
$f^{-1}(U)$ est intègre, cela entraînera $g(U)=f^{-1}(U)$, d'où la
conclusion.

Considérons alors deux cas :

\textup{I)} Supposons d'abord $Y$ normal, et soit $V$ l'ensemble des
points de $X$ en lesquels $f$ est quasi-fini; c'est un ensemble ouvert
dans $X$ {\hyperlink{ega4.c04.cor.13-1-4}{(13.1.4)}}.\footnote{Le
lecteur vérifiera que ni {\hyperlink{ega4.c04.cor.8-12-12}{(8.12.12)}}
ni aucune de ses conséquences n'est utilisée dans la preuve de
{\hyperlink{ega4.c04.cor.13-1-4}{(13.1.4)}}.} Il résulte alors de
{\hyperlink{ega4.c04.cor.8-12-10}{(8.12.10)}}

\oldpage[IV]{353}
que la restriction $f|V:V\to Y$ est une immersion ouverte; donc il y a
un voisinage ouvert $U$ de $y\in Y$ tel que $f|V$ soit un isomorphisme
de $V$ sur le sous-préschéma induit par $Y$ sur $U$; l'isomorphisme
réciproque $g$ est alors une $U$-section de $f^{-1}(U)$ telle que
$g(y)=x$.

\textup{II)} L'obtention d'une $U$-section $g$ telle que $g(y)=x$ étant
un problème local sur $X$ et sur $Y$, on peut supposer que $X$ et $Y$
sont affines et que $f$ est de type fini, donc que $X$ est un
sous-préschéma fermé d'un $Y$-préschéma $Z$ de présentation finie sur
$Y$ (par exemple de la forme $Y[T_1,\ldots,T_n]$). Posons
\[
 Y'=\operatorname{Spec}(\sh{O}_{Y,y}),\qquad
 X'=X_{(Y')},\qquad Z'=Z_{(Y')},
\]
de sorte que $X'$ est un sous-préschéma fermé de $Z'$. En outre, si
$f'=f_{(Y')}$, on a $f'^{-1}(y)=f^{-1}(y)$ (I, 3.6.5) et $x$ est donc
isolé dans $f'^{-1}(y)$. Puisque $Y'$ est normal, on peut appliquer le
résultat de I) à $f'$, et puisque $Y'$ est le seul voisinage de $y$ dans
$Y'$, il existe une $Y'$-section $g'$ de $X'$ telle que $g'(y)=x$, que
l'on peut évidemment aussi considérer comme une $Y'$-section de $Z'$.
Puisque le morphisme structural $h:Z\to Y$ est de présentation finie,
on peut appliquer ({\hyperlink{ega4.c04.thm.8-8-2}{8.8.2}}, (i) et
(ii)) suivant la méthode de
({\hyperlink{ega4.c04.statement.8-1-2}{8.1.2}}, a)), d'où l'on conclut
à l'existence d'un voisinage ouvert $U$ de $y$ dans $Y$ et d'une
$U$-section $g$ de $h^{-1}(U)$ telle que $g(y)=x$. Il reste à voir
qu'en fait $g$ est une $U$-section de $f^{-1}(U)$. Mais puisque
$f^{-1}(U)$ est un sous-préschéma fermé de $h^{-1}(U)$ et que $U$ est
réduit, il suffit pour cela (I, 5.2.2) de voir que
$g(U)\subset f^{-1}(U)$. Or, puisque $U$ est intègre et $f^{-1}(U)$
fermé dans $h^{-1}(U)$, il suffit de prouver que si $\eta$ est le point
générique de $Y$ (et de $U$), on a $g(\eta)\in f^{-1}(U)$. Or, cela
résulte de ce que la restriction de $g$ à $Y'$ est une $Y'$-section de
$X'\subset f^{-1}(U)$ et de ce que $\eta\in Y'$.
\end{proof}

\paragraph{(Err\textsubscript{IV}, 31).}
Dans ({\hyperlink{ega4.c04.subsec.10-4}{IV, 10.4}}), ajouter, après la
proposition {\hyperlink{ega4.c04.statement.10-4-11}{(10.4.11)}} :

\begin{corollary}[10.4.12]
\label{IV.10.4.12}
Soient $S$ un préschéma, $X$ un $S$-préschéma de présentation finie,
$f$ un $S$-endomorphisme radiciel de $X$.
\begin{enumerate}[label=\textup{(\roman*)}]
\item L'endomorphisme $f$ est bijectif et fini (donc un homéomorphisme
universel {\hyperlink{ega4.c04.prop.2-4-5}{(2.4.5)}}).
\item Supposons de plus que $X$ soit réduit, que les anneaux locaux de
$S$ soient noethériens et universellement japonais, et que $f$ soit
birationnel (condition automatiquement vérifiée si pour tout point
maximal $x$ de $X$, $k(x)$ est un corps parfait). Alors $f$ est un
$S$-automorphisme.
\end{enumerate}
\end{corollary}

\begin{proof}
\textup{(i)} La question étant locale sur $S$, on peut se borner au cas
où $S=\operatorname{Spec}(A)$ est affine. Il existe alors une
sous-$\mathbb Z$-algèbre de type fini $A_0$ de $A$ telle qu'en posant
$S_0=\operatorname{Spec}(A_0)$, on ait
$X=X_0\times_{S_0}S$, où $X_0$ est un préschéma de type fini sur $S_0$,
et que $f$ se déduise d'un $S_0$-endomorphisme $f_0$ de $X_0$ par le
changement de base $S\to S_0$
{\hyperlink{ega4.c04.prop.8-9-1}{(8.9.1)}}; on peut en outre supposer
$f_0$ radiciel ({\hyperlink{ega4.c04.thm.8-10-5}{8.10.5}}, (vii)). Il
suffira de prouver que $f_0$ est bijectif et fini, et on peut donc
supposer désormais que $A$ est une $\mathbb Z$-algèbre de type fini, ce
qui entraîne que $S$ et $X$ sont des préschémas excellents
({\hyperlink{ega4.c04.statement.7-8-3}{(7.8.3)}} et
{\hyperlink{ega4.c04.prop.7-8-6}{(7.8.6)}}) et en particulier que $X$
est un préschéma universellement japonais
({\hyperlink{ega4.c04.cor.5-11-4}{(5.11.4)}} et
({\hyperlink{ega4.c04.statement.7-8-3}{7.8.3}}, (vi))). On peut en outre
supposer $X$ réduit; en effet,
$f_{\mathrm{red}}:X_{\mathrm{red}}\to X_{\mathrm{red}}$ est un
$S_{\mathrm{red}}$-endomorphisme radiciel (I, 5.1.6); si l'on prouve que
$f_{\mathrm{red}}$ est fini, il en sera de même de $f$ : en effet,
puisque $f$ est radiciel, il est séparé
{\hyperlink{ega4.c04.lem.1-8-7-1}{(1.8.7.1)}}; comme
$f_{\mathrm{red}}$ est propre il en est de même de $f$ (II, 5.4.6);
mais puisque $f$ est évidemment quasi-fini, il est fini (III, 4.4.2).

\oldpage[IV]{354}
Soit alors $X'$ le normalisé de $X$ (II, 6.3.8); si $X_i$
($1\leq i\leq n$) sont les sous-préschémas réduits de $X$ ayant pour
espaces sous-jacents les composantes irréductibles de $X$, $X'$ est
somme des normalisés $X'_i$ des préschémas intègres $X_i$; nous
désignerons par $x_i$ (resp. $x'_i$) le point générique de $X_i$
(resp. $X'_i$). Le morphisme canonique $p:X'\to X$ est fini puisque
$X$ est universellement japonais et on a
$p^{-1}(x_i)=\{x'_i\}$, l'homomorphisme $k(x_i)\to k(x'_i)$ étant
bijectif.

Soit $f=(\psi,\theta)$; on sait déjà
{\hyperlink{ega4.c04.statement.10-4-11}{(10.4.11)}} que $\psi$ est une
bijection continue de l'espace $X$ sur lui-même; si l'on pose
$y_i=\psi(x_i)$, $y_i$ est aussi un point maximal de $X$, car si $y'_i$
était une générisation de $y_i$ distincte de $y_i$, $x_i$ serait une
générisation $\neq x_i$ de $\psi^{-1}(y'_i)$. Puisque l'ensemble $M$
des points maximaux de $X$ est fini, la restriction $\psi|M$ est une
bijection de $M$ sur lui-même, et on peut écrire
$\psi(x_i)=x_{\sigma(i)}$, où $\sigma$ est une permutation de
$[1,n]$. Alors, pour tout $i$, la restriction $f|X_i$ de $f$ se
factorise en
\[
 X_i\xrightarrow{f_i}X_{\sigma(i)}\longrightarrow X
\]
(I, 5.2.2), où la seconde flèche est l'injection canonique. Comme les
$X_i$ sont intègres, il existe (II, 6.3.9) un unique morphisme
$f'_i:X'_i\to X'_{\sigma(i)}$ rendant commutatif le diagramme
\[
\begin{array}{ccc}
X'_i&\xrightarrow{f'_i}&X'_{\sigma(i)}\\
\downarrow p_i&&\downarrow p_{\sigma(i)}\\
X_i&\xrightarrow{f_i}&X_{\sigma(i)}
\end{array}
\]
et si $f'$ coïncide avec $f'_i$ dans chacun des $X'_i$, on a donc un
diagramme commutatif
\phantomsection\label{IV.10.4.12.1}
\begin{equation}
\label{IV.10.4.12.1.errata}
\begin{array}{ccc}
X'&\xrightarrow{f'}&X'\\
\downarrow p&&\downarrow p\\
X&\xrightarrow{f}&X.
\end{array}
\tag{10.4.12.1}
\end{equation}

D'autre part, comme $f$ est séparé, il en est de même de
$f\circ p=p\circ f'$, donc $f'$ est séparé (I, 5.5.1); enfin, puisque
$p$ est fini et $f$ bijectif, il résulte de
{\hyperlink{ega4.c04.sec.10-4-12-1}{(10.4.12.1)}} que $f'$ est
quasi-fini. Puisque les $X'_i$ sont intègres et normaux, il résulte de
{\hyperlink{ega4.c04.cor.8-12-11}{(8.12.11)}} que $f'_i$ se factorise en
\[
 X'_i\xrightarrow{f''_i}X''_{\sigma(i)}
 \xrightarrow{u_{\sigma(i)}}X'_{\sigma(i)},
\]
où $X''_{\sigma(i)}$ est la fermeture intégrale de
$X'_{\sigma(i)}$ relative à l'extension $k(x'_i)$ de
$k(x'_{\sigma(i)})$, et $f''_i$ une immersion ouverte. Mais comme
$k(x'_i)$ est une extension radicielle de $k(x'_{\sigma(i)})$,
$u_{\sigma(i)}$ est un morphisme radiciel (Bourbaki,
\emph{Alg. comm.}, chap.~V, \S2, n\textsuperscript{o}~3, lemme 4);
donc $f'$ est radiciel, et il résulte alors de
{\hyperlink{ega4.c04.statement.10-4-11}{(10.4.11)}} que $f'$ est
surjectif; on en conclut (puisque les $u_{\sigma(i)}$ sont injectifs)
que $f''_i$ est bijectif, donc un isomorphisme; on en déduit enfin que
le morphisme $f'$ est fini (puisqu'il en est ainsi de $u_{\sigma(i)}$,
$X'$ étant universellement japonais
{\hyperlink{ega4.c04.statement.7-8-3}{(7.8.3)}}). Cela étant, comme
$p\circ f'=f\circ p$ est alors fini et que $p$ est surjectif et $f$
séparé, $f$ est propre (II, 5.4.3) et quasi-fini, donc fini
{\hyperlink{ega4.c04.thm.8-11-1}{(8.11.1)}}.

\oldpage[IV]{355}
\textup{(ii)} Compte tenu de
({\hyperlink{ega4.c04.thm.8-10-5}{8.10.5}}, (i)) appliqué suivant la
méthode de ({\hyperlink{ega4.c04.statement.8-1-2}{8.1.2}}, a)), on peut
se borner au cas où $S$ est un schéma local. Les réductions préalables
faites dans la preuve de (i) sont donc ici inutiles, et avec les mêmes
notations on voit que l'on a un diagramme commutatif
{\hyperlink{ega4.c04.sec.10-4-12-1}{(10.4.12.1)}} dans lequel $X$ est
un $S$-préschéma de type fini, $p$ est fini, $f'$ est un isomorphisme et
l'homomorphisme canonique $\sh{O}_X\to p_*(\sh{O}_{X'})$ est injectif.
Mais on montrera au chap.~VI, dans la théorie des quotients des
préschémas par des relations d'équivalence, que sous ces conditions $f$
est lui-même un isomorphisme (le lecteur vérifiera qu'on n'utilisera pas
{\hyperlink{ega4.c04.cor.10-4-12}{(10.4.12)}} dans la théorie des
relations d'équivalence du chap.~VI).
\end{proof}

\begin{corollary}[10.4.13]
\label{IV.10.4.13}
Soient $k$ un corps, $X$ un préschéma algébrique sur $k$. Alors tout
$k$-endomorphisme radiciel $f$ de $X$ est fini et bijectif, donc un
homéomorphisme universel. Si de plus $X$ est réduit et $f$ birationnel
(par exemple si $k$ est de caractéristique 0), $f$ est un
$k$-automorphisme.
\end{corollary}

\begin{remark}[10.4.14]
\label{IV.10.4.14}
\begin{enumerate}[label=\textup{(\roman*)}]
\item Dans {\hyperlink{ega4.c04.cor.10-4-12}{(10.4.12)}} et
{\hyperlink{ega4.c04.cor.10-4-13}{(10.4.13)}} on ne peut supprimer
l'hypothèse que $X$ est réduit pour la seconde conclusion; c'est ce que
montre l'exemple du $B$-endomorphisme correspondant à l'augmentation
$D_B(L)\to B$ de l'anneau $D_B(L)=A$ défini dans ($0$, 18.2.3), où
$B$ est identifié à un sous-anneau de $A$.

\item Soit $X$ un préschéma noethérien. Nous ignorons si tout
endomorphisme radiciel $f$ de $X$ est nécessairement surjectif et fini.
Si $f$ est supposé surjectif et si de plus $X$ est un préschéma
japonais, le raisonnement de
{\hyperlink{ega4.c04.cor.10-4-12}{(10.4.12)}} montre que $f$ est fini;
si de plus $X$ est réduit et si $f$ est birationnel, par exemple si,
pour tout point maximal $x\in X$, $k(x)$ est parfait, le même
raisonnement montre que $f$ est un automorphisme.

\item On peut donner des exemples d'endomorphismes radiciels non
surjectifs de schémas affines intègres non noethériens : si $k$ est un
anneau, $A=k[(T_\lambda)_{\lambda\in L}]$ un anneau de polynômes où $L$
est infini dénombrable, et $\mathfrak J$ l'idéal engendré dans $A$ par
une sous-famille $(T_\lambda)_{\lambda\in H}$, où $L-H$ est infini et
$H\neq\varnothing$, on remarque que $A/\mathfrak J$ est isomorphe à
$k[(T_\lambda)_{\lambda\in L-H}]$, donc à $A$; l'endomorphisme de
$\operatorname{Spec}(A)$ correspondant à l'endomorphisme d'anneaux
composé
\[
 A\xrightarrow{p}A/\mathfrak J\xrightarrow{u}A
\]
où $p$ est la projection canonique et $u$ un isomorphisme, répond à la
question, son image étant le sous-préschéma fermé de
$\operatorname{Spec}(A)$ défini par $\mathfrak J$.
\end{enumerate}
\end{remark}

\paragraph{(Err\textsubscript{IV}, 32).}
Dans (IV, 11.3.14), ligne 7 de la p.~141, avant « Soient », insérer :
(i). Après la ligne 9 de la p.~141, ajouter :

\smallskip
\noindent
\textup{(ii)} Inversement, soit $f:X\to Y$ un morphisme plat. Si $X$
est unibranche (resp. géométriquement unibranche) en un point $x$, $Y$
est unibranche (resp. géométriquement unibranche) au point $y=f(x)$.

Au début de la ligne 10 de la p.~141, insérer : (i). Après la ligne 20
de la p.~141, ajouter :

\smallskip
\noindent
\textup{(ii)} Notons que si $Y_1=Y_{\mathrm{red}}$,
$X_1=X\times_Y Y_1$ a pour espace sous-jacent celui de $X$ et le
morphisme $f_1:X_1\to Y_1$ déduit de $f$ par changement de base est
plat. On peut donc se borner au cas où $Y=\operatorname{Spec}(A)$,
$A$ étant un anneau local réduit, et $y$ le point fermé de $Y$. Posons
$B=\sh{O}_{X,x}$; par hypothèse $B_{\mathrm{red}}$ est unibranche, donc
$\operatorname{Spec}(B)$ n'a qu'un seul point maximal, et par platitude
{\hyperlink{ega4.c04.prop.2-3-4}{(2.3.4)}} on en déduit qu'il en est de
même

\oldpage[IV]{356}
de $\operatorname{Spec}(A)$; puisque $A$ est réduit, il est donc
intègre. Soient $K$ le corps des fractions de $A$, $A'$ la clôture
intégrale de $A$; par platitude, $B\otimes_AA'$ s'identifie à un
sous-anneau de $B\otimes_AK$ contenant $B$, et comme les éléments non
nuls de $A$ sont $B$-réguliers par platitude, $B\otimes_AK$ s'identifie
à un sous-anneau de l'anneau total des fractions $R$ de $B$. Pour
montrer que $A$ est unibranche, il suffit de voir que $A'$ ne peut avoir
deux idéaux maximaux distincts (nécessairement au-dessus de $y$
(Bourbaki, \emph{Alg. comm.}, chap.~V, \S2,
n\textsuperscript{o}~1, prop.~1)); s'il en était ainsi,
$B'=B\otimes_AA'$ aurait deux idéaux premiers distincts au-dessus de
$x$ (I, 3.4.7), et il suffit de voir que ce n'est pas possible.

Soit $\mathfrak N$ le nilradical de $B$; si $S$ est l'ensemble des
éléments réguliers de $B$, on a $R=S^{-1}B$, et
$S^{-1}\mathfrak N=R\mathfrak N$ est le nilradical de $R$, donc
\[
 R/R\mathfrak N=S^{-1}B/S^{-1}\mathfrak N=S^{-1}(B/\mathfrak N)
\]
est contenu dans le corps des fractions $L$ de l'anneau local intègre
$B/\mathfrak N=B_{\mathrm{red}}$; il en est donc de même de
\[
 B'_{\mathrm{red}}=B'/(B'\cap R\mathfrak N)
\]
qui est entier sur $B_{\mathrm{red}}$; par hypothèse
$B'_{\mathrm{red}}$, étant contenu dans la clôture intégrale de
$B_{\mathrm{red}}$, est un anneau local, donc il en est de même de
$B'$, ce qui établit notre assertion.

Si l'on suppose en outre $B$ géométriquement unibranche, le corps
résiduel de $B'_{\mathrm{red}}$, étant contenu dans le corps résiduel de
la clôture intégrale de $B_{\mathrm{red}}$, est une extension radicielle
du corps résiduel $k(x)$ de $B$. Notons maintenant que le corps résiduel
de $B'_{\mathrm{red}}$ est aussi celui de $B'=B\otimes_AA'$, donc est
égal à $k(x)\otimes_{k(y)}k'$, où $k'$ est le corps résiduel de $A'$;
on déduit alors de {\hyperlink{ega4.c04.prop.2-6-1}{(2.6.1)}} que $k'$
est une extension radicielle de $k(y)$, donc $A$ est aussi
géométriquement unibranche.

\paragraph{(Err\textsubscript{IV}, 33).}
Dans ({\hyperlink{ega4.c04.item.12-1-2}{IV, 12.1.2}}), l'exemple donné
en (i) n'est pas correct, le morphisme considéré étant ouvert. On
construit un exemple d'un morphisme plat $f:X\to Y$ dont la restriction
à une composante irréductible de $X$ n'est pas un morphisme ouvert de
la façon suivante. Soit $Z=\operatorname{Spec}(\mathbb C[U,V])$ le
plan affine sur le corps des nombres complexes, et considérons le
schéma $Y$ obtenu en recollant deux points fermés distincts $a$, $b$ de
$Z$; on a $Y=\operatorname{Spec}(B)$, où $B$ est le sous-anneau de
$\mathbb C[U,V]$ formé des polynômes prenant la même valeur aux points
$a$ et $b$ (l'ensemble des points fermés de $Z$ étant identifié à
$\mathbb C^2$); il est immédiat que $B$ est un sous-$\mathbb C$-espace
vectoriel de $\mathbb C[U,V]$ de codimension 1, donc $\mathbb C[U,V]$
est un $B$-module de type fini, et par suite le morphisme $Z\to Y$ est
fini, l'image réciproque du point $c$ de $Y$ obtenu par recollement de
$a$ et de $b$ étant l'ensemble $\{a,b\}$.

Considérons alors le schéma $X$ obtenu en recollant deux exemplaires de
$Z$, soient $Z_1$, $Z_2$, de sorte que si $a_i$, $b_i$ ($i=1,2$) sont
les points de $Z_i$ correspondant à $a$ et $b$, $a_1$ soit recollé à
$b_2$ et $b_1$ à $a_2$ (cf. ({\hyperlink{ega4.c04.item.6-5-5}{6.5.5}},
(ii))). On voit aisément (\emph{loc. cit.}) que le morphisme $X\to Y$
ainsi obtenu est plat et fini (et même étale) et que $Z_1$ et $Z_2$
s'identifient aux deux composantes irréductibles de $X$. Cependant, si
$D$ est une droite dans $Z_1$ passant par $a_1$ mais non par $b_1$,
l'image $V$ dans $Y$ de l'ouvert $Z_1-D$ de $Z_1$ contient l'image de
$b_1$, c'est-à-dire le point $c$, mais ne contient pas l'image du point
générique de $D$, donc $V$ n'est pas un voisinage de $c$.

\paragraph{(Err\textsubscript{IV}, 34).}
Dans le titre de (IV, 13.2), ligne 5 du bas de la p.~190, remplacer :
« équidimensionnels » par « localement équidimensionnels ». De même
dans (IV, 13.2.2), ligne 4 de la p.~191, remplacer deux fois
« équidimensionnel » par « localement équidimensionnel ».

\oldpage[IV]{357}
Ligne 5 de la p.~191, ajouter : \emph{Si de plus les fibres
$f^{-1}(f(x))$ de $f$ sont équidimensionnelles ($0_{\mathrm{IV}}$,
14.1.3), on dit que $f$ est équidimensionnel.} Ligne 6 du bas de la
p.~194, remplacer « équidimensionnelles » par « localement
équidimensionnelles ».

\paragraph{(Err\textsubscript{IV}, 35).}
Dans le titre de (IV, 13.3), ligne 4 du bas de la p.~94, remplacer
« équidimensionnels » par « localement équidimensionnels ». Dans
(IV, 13.3.2), ligne 6 du bas de la p.~196, remplacer deux fois
« équidimensionnel » par « localement équidimensionnel ». Ligne 5 du bas
de la p.~196, ajouter : \emph{Si de plus les fibres $f^{-1}(f(x))$ de
$f$ sont équidimensionnelles ($0_{\mathrm{IV}}$, 14.1.3), on dit que
$f$ est équidimensionnel.} Dans (IV, 13.3.7), ligne 14 de la p.~198,
remplacer « équidimensionnel » par « localement équidimensionnel ».
Dans ({\hyperlink{ega4.c04.statement.13-3-9}{IV, 13.3.9}}), ligne 3 du
bas de la p.~198, remplacer « équidimensionalité » par
« équidimensionalité locale ». Lignes 10, 12, 20 et 24 de la p.~199,
remplacer « équidimensionnel » par « localement équidimensionnel ».

\paragraph{(Err\textsubscript{IV}, 36).}
Dans l'introduction du \S14, ligne 6 du bas de la p.~199, remplacer
« équidimensionnel » par « localement équidimensionnel ».

\paragraph{(Err\textsubscript{IV}, 37).}
Dans (IV, 14.4.1), ligne 5 du bas de la p.~209, remplacer « localement
de type fini » par « localement de présentation finie ». L'énoncé
actuel de ({\hyperlink{ega4.c04.thm.14-4-1}{IV, 14.4.1}}) est inexact,
même lorsque $f:X\to Y$ est une immersion fermée, $X$ étant réduit à
un seul point $y$ (fermé) de $Y$, et $\sh{O}_{Y,y}$ étant un corps. Il
y a en effet des anneaux $A$ tels que $Y=\operatorname{Spec}(A)$ soit
compact, sans point isolé et totalement discontinu, tout idéal premier
$\mathfrak p$ de $A$ étant maximal et $A_\mathfrak p$ étant un corps
(Bourbaki, \emph{Alg. comm.}, chap.~II, \S4, exerc.~16 et 17); il est
clair que l'immersion fermée $\{y\}\to Y$ n'est pas un morphisme
ouvert, mais vérifie les conditions b) et b') de
({\hyperlink{ega4.c04.thm.14-4-1}{IV, 14.4.1}}).

\paragraph{(Err\textsubscript{IV}, 38).}
Dans (IV, 14.4.1.3), ligne 14 de la p.~211, remplacer « localement
quasi-fini et dominant » par « localement quasi-fini, localement de
présentation finie et dominant ».

\paragraph{(Err\textsubscript{IV}, 39).}
Dans (IV, 14.4.2), ligne 10 du bas de la p.~211, remplacer « localement
de type fini » par « localement de présentation finie ».

\paragraph{(Err\textsubscript{IV}, 40).}
Dans (IV, 14.4.3), ligne 18 de la p.~212, remplacer « localement de
type fini » par « localement de présentation finie ».

\paragraph{(Err\textsubscript{IV}, 41).}
Dans ({\hyperlink{ega4.c04.cor.14-4-4}{IV, 14.4.4}}), lignes 15 et 16 du
bas de la p.~212, remplacer « localement de type fini » par
« localement de présentation finie ». Lignes 10 et 11 du bas, remplacer
« équidimensionnel » par « localement équidimensionnel ».

\paragraph{(Err\textsubscript{IV}, 42).}
Dans (IV, 14.4.8), après la ligne 1 de la p.~214, ajouter :
\begin{quote}
\emph{Si de plus l'ensemble des composantes irréductibles de $Y$ est
localement fini, ces deux conditions sont équivalentes aux suivantes :}
\end{quote}

\paragraph{(Err\textsubscript{IV}, 43).}
Dans (IV, 14.4.10), ligne 15 de la p.~215, remplacer : quelconque, par :
dont l'ensemble des composantes irréductibles est localement fini.
Ligne 16 de la p.~215, remplacer : de type fini, par : de présentation
finie.

\paragraph{(Err\textsubscript{IV}, 44).}
Dans (IV, 15.5.6.1), ligne 15 du bas de la p.~234, après la première
phrase, ajouter :

\begin{quote}
On peut se limiter au cas où dans $X$ toute partie fermée irréductible
n'admet qu'un seul point générique. En effet, dans le cas contraire on
considère l'espace $Y=X/R$,

\oldpage[IV]{358}
quotient de $X$ par la relation d'équivalence
$\overline{\{x\}}=\overline{\{x'\}}$. Il résulte aussitôt de cette
définition que tout ouvert de $X$ est saturé pour $R$, donc la topologie
de $X$ est l'image réciproque de celle de $Y$ par l'application
canonique $\pi:X\to X/R$. Un ouvert dense dans $X$ est de la forme
$\pi^{-1}(U)$, où $U$ est un ouvert dense dans $Y$, donc un fermé rare
$Z$ dans $X$ est de la forme $\pi^{-1}(S)$, où $S=\pi(Z)$ est fermé
rare dans $Y$, et il revient au même de dire que $X-Z$ est connexe ou
que $Y-S$ est connexe. De même $\pi(T_x)$ est l'ensemble des
générisations de $\pi(x)$, et il est équivalent de dire que
$T_x-\{x\}$ est connexe ou que $\pi(T_x)-\{\pi(x)\}$ l'est. Enfin, par
construction $Y$ est un espace de Kolmogoroff localement noethérien,
dont toute partie fermée admet donc un seul point générique. On voit
alors qu'il revient au même de prouver le lemme pour $X$ ou pour $Y$,
et on peut donc supposer que toute partie fermée irréductible de $X$ a
un point générique unique.
\end{quote}

\paragraph{(Err\textsubscript{IV}, 45).}
Dans (IV, 15.6.1), ligne 18 du bas de la p.~236, supprimer « localement
noethérien » (hypothèse non utilisée dans la démonstration).

\paragraph{(Err\textsubscript{IV}, 46).}
Dans (IV, 15.6.3), ligne 11 de la p.~237, au lieu de « supposons que
pour tout $y\in Y$ », lire : « supposons que $f$ soit localement de
présentation finie et que pour tout $y\in Y$ ». Remplacer la
démonstration (lignes 14 à 24 de la p.~237) par la suivante :
\begin{quote}
En vertu de {\hyperlink{ega4.c04.prop.1-10-3}{(1.10.3)}}, il suffit de
prouver que pour tout $x\in X_y^0$ et toute générisation $y'$ de $y$,
il existe une générisation $x'$ de $x$ telle que $f(x')=y'$. Or il
suffit de prendre pour $x'$ le point générique de $X_{y'}^0$, en vertu
de {\hyperlink{ega4.c04.cor.15-6-2}{(15.6.2)}}.
\end{quote}

\paragraph{(Err\textsubscript{IV}, 47).}
Dans (IV, 15.6.7), ligne 2 de la p.~241, après « noethérien », ajouter :
« ou si la condition $\beta$) de
({\hyperlink{ega4.c04.prop.15-6-6}{15.6.6}}, (iii)) est vérifiée ».
Après la ligne 9 de la p.~241, ajouter :
\begin{quote}
Lorsque la condition $\beta$) de
({\hyperlink{ega4.c04.prop.15-6-6}{15.6.6}}, (iii)) est satisfaite, le
fait que a) entraîne b') résulte de
({\hyperlink{ega4.c04.prop.15-6-6}{15.6.6}}, (ii)), et le fait que b')
entraîne a) résulte de
({\hyperlink{ega4.c04.prop.15-6-6}{15.6.6}}, (iii)).
\end{quote}

\paragraph{(Err\textsubscript{IV}, 48).}
Dans (IV, 20.1.5), ligne 13 de la p.~228, (IV, 20.1.6), ligne 17 de la
p.~228, (IV, 20.1.7), ligne 13 du bas de la p.~228, supprimer
« strictement ».

\paragraph{(Err\textsubscript{IV}, 49).}
Dans (IV, 20.2.13), ajouter après la ligne 6 de la p.~236 : En
particulier si $X$ est réduit et n'a qu'un nombre fini de composantes
irréductibles, la notion de $\sh{O}_X$-Module sans torsion définie dans
({\hyperlink{ega4.c04.statement.20-1-5}{IV, 20.1.5}}) coïncide avec
celle de (I, 7.4.1).

\paragraph{(Err\textsubscript{IV}, 50).}
Dans (IV, 20.2.14), lignes 18 et 25 de la p.~236, supprimer
« strictement ».

\paragraph{(Err\textsubscript{IV}, 51).}
Dans (IV, 20.6.2), ligne 13 du bas de la p.~252, supprimer
« strictement ».

\paragraph{(Err\textsubscript{IV}, 52).}
Dans (IV, 20.6.4), lignes 23 et 22 de la p.~253, supprimer
« strictement ».

\paragraph{(Err\textsubscript{IV}, 53).}
Dans (IV, 21.4), après la ligne 23 de la p.~267, ajouter :

La proposition suivante et ses corollaires nous ont été signalés par
M.~Raynaud.

\begin{proposition}[21.4.9]
\label{IV.21.4.9}
Soient $X$, $Y$ deux préschémas localement noethériens,
$f:X\to Y$ un morphisme fidèlement plat; on suppose en outre, soit que
$f$ est quasi-compact, soit que $f$ est ouvert. Pour qu'un diviseur
$D'$ sur $X$ soit de la forme $f^*(D)$, où $D$ est un diviseur sur
$Y$, il faut et il suffit que pour tout $y\in Y$ tel que
$\operatorname{prof}(\sh{O}_{Y,y})\leq1$, en posant
\[
 Y_1=\operatorname{Spec}(\sh{O}_{Y,y}),\qquad
 X_1=X\times_YY_1,\qquad
 f_1=f_{(Y_1)}:X_1\to Y_1,
\]
l'image réciproque $D'_1$ de $D'$ sur $X_1$ soit de la forme
$f_1^*(D_1)$, où $D_1\in\operatorname{Div}(Y_1)$.
\end{proposition}

\oldpage[IV]{359}
\begin{proof}
La nécessité de la condition est évidente par transitivité des images
réciproques de diviseurs, tous les morphismes considérés étant plats
{\hyperlink{ega4.c04.statement.21-4-4}{(21.4.4)}}. Pour tout ouvert
$U$ de $Y$, notons $f_U$ la restriction $f^{-1}(U)\to U$ de $f$; en
vertu de {\hyperlink{ega4.c04.prop.21-4-8}{(21.4.8)}} l'application
\[
 D_U\longmapsto f_U^*(D_U)
\]
de $\operatorname{Div}(U)$ dans $\operatorname{Div}(f^{-1}(U))$ est
injective; donc, si $U_1$, $U_2$ sont deux ouverts de $Y$, et si
\[
 D'|f^{-1}(U_1)=f_{U_1}^*(D_1),\qquad
 D'|f^{-1}(U_2)=f_{U_2}^*(D_2),
\]
on a nécessairement
$D_1|U_1\cap U_2=D_2|U_1\cap U_2$. Ceci montre aussitôt qu'il existe
un plus grand ouvert $U\subset Y$ tel que
\[
 D'|f^{-1}(U)=f_U^*(D_U),
 \qquad D_U\in\operatorname{Div}(U);
\]
il s'agit de voir que $U=Y$. Supposons le contraire, et soit $y$ un
point maximal de $Y-U$; on va montrer qu'il existe un voisinage ouvert
$V$ de $y$ dans $Y$ tel que
$D'|f^{-1}(V)=f_V^*(D_V)$ pour un $D_V\in\operatorname{Div}(V)$, ce qui
contredira le caractère maximal de $U$.

\textup{I) Cas où $f$ est quasi-compact.} Appliquant le fait que $f$
est quasi-compact et {\hyperlink{ega4.c04.prop.20-3-8}{(20.3.8)}}
suivant la méthode de ({\hyperlink{ega4.c04.statement.8-1-2}{8.1.2}},
a)), on est ramené à prouver qu'avec les notations de l'énoncé (mais
sans supposer d'abord de limitation sur
$\operatorname{prof}(\sh{O}_{Y,y})$), $D'_1$ est de la forme
$f_1^*(D_1)$. Or, si $\operatorname{prof}(\sh{O}_{Y,y})\leq1$, cette
assertion n'est autre que l'hypothèse. Supposons donc
$\operatorname{prof}(\sh{O}_{Y,y})\geq2$; puisque $y$ est point maximal
de $Y-U$, l'ouvert $W=Y_1\cap U$ de $Y_1$ est égal à $Y_1-\{y\}$, et
on peut se borner au cas où $Y=Y_1$, $U=Y-\{y\}$. Par hypothèse, il
existe un diviseur $D_U\in\operatorname{Div}(U)$ tel que
$f_U^*(D_U)=D'|f^{-1}(U)$; $D_U$ est classe d'équivalence d'un couple
$(\sh{L}_U,s)$, où $\sh{L}_U$ est un $\sh{O}_U$-Module inversible et
$s$ une section méromorphe régulière de $\sh{L}_U$ au-dessus de $U$
{\hyperlink{ega4.c04.prop.21-2-11}{(21.2.11)}}; par suite
{\hyperlink{ega4.c04.statement.21-4-1}{(21.4.1)}},
$D'|f^{-1}(U)$ est classe d'équivalence du couple
$(\sh{L}'_U,s')$, où
\[
 \sh{L}'_U=f_U^*(\sh{L}_U),\qquad s'=s\circ f_U.
\]
Notons $i:U\to Y$, $j:f^{-1}(U)\to X$ les injections canoniques, et
posons
\[
 \sh{L}=i_*(\sh{L}_U),\qquad
 \sh{L}'=j_*(\sh{L}'_U).
\]
Comme $f$ est plat, $f^*(\sh{L})$ est canoniquement isomorphe à
$\sh{L}'$ {\hyperlink{ega4.c04.lem.2-3-1}{(2.3.1)}}; d'autre part,
puisque $f$ est plat, la relation
$\operatorname{prof}(\sh{O}_{Y,y})\geq2$ entraîne
$\operatorname{prof}(\sh{O}_{X,x})\geq2$ pour tout $x\in f^{-1}(y)$
{\hyperlink{ega4.c04.prop.6-3-1}{(6.3.1)}}.

Or, $D'$ est classe d'équivalence d'un couple $(\sh{L}'',s'')$, où
$\sh{L}''$ est un $\sh{O}_X$-Module inversible et $s''$ une section
méromorphe régulière de $\sh{L}''$ au-dessus de $X$; $\sh{L}'_U$ est
par suite isomorphe à $\sh{L}''|f^{-1}(U)$, et il résulte alors de ce
qui précède et de {\hyperlink{ega4.c04.prop.5-9-8}{(5.9.8)}} et
{\hyperlink{ega4.c04.thm.5-10-5}{(5.10.5)}} que $\sh{L}'$ est un
$\sh{O}_X$-Module inversible isomorphe à $\sh{L}''$. Puisque $f$ est
fidèlement plat et quasi-compact, on en conclut que $\sh{L}$ est un
$\sh{O}_Y$-Module inversible
{\hyperlink{ega4.c04.prop.2-5-2}{(2.5.2)}} et par suite isomorphe à
$\sh{O}_Y$, puisque $Y$ est un schéma local. D'ailleurs, puisque
$\operatorname{prof}(\sh{O}_{Y,y})\geq1$, $U=Y-\{y\}$ est
schématiquement dense dans $Y$
({\hyperlink{ega4.c04.statement.20-2-13}{20.2.13}}, (iv)) et par suite
la section méromorphe régulière $s$ de $\sh{L}_U$ au-dessus de $U$ se
prolonge en une section méromorphe régulière $s_0$ de $\sh{L}$
au-dessus de $Y$ {\hyperlink{ega4.c04.prop.20-2-11}{(20.2.11)}}; la
classe d'équivalence de $(\sh{L},s_0)$ est donc un diviseur $D$ sur
$Y$; et il est clair que $f^*(D)$ et $D'$ induisent le même diviseur sur
$f^{-1}(U)=X-f^{-1}(y)$; mais comme
$\operatorname{prof}(\sh{O}_{X,x})\geq2$ pour tout $x\in f^{-1}(y)$,
on a $D'=f^*(D)$ {\hyperlink{ega4.c04.prop.21-1-8}{(21.1.8)}}, ce qui
achève la démonstration dans ce cas.

\textup{II) Cas où $f$ est ouvert.} Comme il s'agit de prouver
l'existence d'un voisinage ouvert $V$ de $y$ dans $Y$ tel que
$D'|f^{-1}(V)=f_V^*(D_V)$ pour un $D_V\in\operatorname{Div}(V)$, on
peut déjà supposer $Y$ affine. Puisque $f$ est ouvert, les images
$f(W)$, où $W$ parcourt les ouverts affines de $X$, forment un
recouvrement ouvert de $Y$, donc, puisque $Y$ est quasi-compact, il
existe un ouvert quasi-compact $T\subset X$ tel que $f(T)=Y$, et le
morphisme $f|T$ est quasi-compact (I, 1.1.1); \emph{a fortiori}, pour
tout ouvert quasi-compact $T_\alpha\supset T$ de $X$,

\oldpage[IV]{360}
$f(T_\alpha)=Y$ et $f|T_\alpha$ est quasi-compact. On peut par suite
appliquer à $f|T_\alpha$ le résultat de I), et il existe un diviseur
$D_\alpha$ sur $Y$ tel que
\[
 (f|T_\alpha)^*(D_\alpha)=D'|T_\alpha;
\]
en vertu de la transitivité des images réciproques de diviseurs par
morphismes plats {\hyperlink{ega4.c04.statement.21-4-4}{(21.4.4)}}, on
a aussi $(f|T)^*(D_\alpha)=D'|T$. On en conclut que pour deux indices
$\alpha$, $\beta$ distincts, on a
\[
 (f|T)^*(D_\alpha)=(f|T)^*(D_\beta),
\]
d'où (par {\hyperlink{ega4.c04.prop.21-4-8}{(21.4.8)}}),
$D_\alpha=D_\beta$. Si $D$ est la valeur commune des $D_\alpha$, on a
donc $f^*(D)|T_\alpha=D'|T_\alpha$ pour tout $\alpha$, et par suite
$f^*(D)=D'$.
\end{proof}

\begin{corollary}[21.4.10]
\label{IV.21.4.10}
Supposons vérifiées les hypothèses générales de
{\hyperlink{ega4.c04.prop.21-4-9}{(21.4.9)}}, et en outre supposons que
$Y$ soit normal et que pour tout $y\in Y$ tel que
$\dim(\sh{O}_{Y,y})=1$, la fibre $X_y=f^{-1}(y)$ n'ait pas de cycle
premier associé immergé. Pour qu'un diviseur $D'$ sur $X$ soit de la
forme $f^*(D)$, où $D\in\operatorname{Div}(Y)$, il faut et il suffit
que les deux conditions suivantes soient vérifiées :
\begin{enumerate}[label=\textup{\arabic*\textdegree}]
\item pour tout point maximal $y$ de $Y$, l'image réciproque de $D'$
dans $\operatorname{Div}(X_y)$ soit nulle;
\item pour tout $y\in Y$ tel que $\dim(\sh{O}_{Y,y})=1$, il existe un
entier $n_y$ tel que, pour tout point maximal $x$ de $X_y$, l'image
réciproque de $D'$ dans
$\operatorname{Div}(\operatorname{Spec}(\sh{O}_{X,x}))$ soit égale à
$\operatorname{div}(t_y^{n_y})$, où $t_y$ est l'image par
l'homomorphisme $\sh{O}_{Y,y}\to\sh{O}_{X,x}$ d'une uniformisante de
l'anneau de valuation discrète $\sh{O}_{Y,y}$.
\end{enumerate}
\end{corollary}

\begin{proof}
Appliquons le critère {\hyperlink{ega4.c04.prop.21-4-9}{(21.4.9)}}.
Aux points maximaux $y$ de $Y$, on a $\operatorname{Div}(Y_1)=0$
puisque $\sh{O}_{Y,y}$ est artinien, et la condition 1\textdegree{} de
(21.4.10) est équivalente au critère de (21.4.9) en un tel point. Il
faut voir que les conditions 1\textdegree{} et 2\textdegree{} impliquent
de même la condition de (21.4.9) en un point $y\in Y$ où
$\operatorname{prof}(\sh{O}_{Y,y})=1$; comme $Y$ est normal, ces points
sont aussi ceux où $\dim(\sh{O}_{Y,y})=1$, ou encore où
$\sh{O}_{Y,y}$ est un anneau de valuation discrète
{\hyperlink{ega4.c04.thm.5-8-6}{(5.8.6)}}. On peut donc se borner au
cas où $Y=\operatorname{Spec}(\sh{O}_{Y,y})$. Si $t$ est une
uniformisante de $\sh{O}_{Y,y}$, l'application
$n\mapsto\operatorname{div}(t^n)$ est un isomorphisme de $\mathbb Z$
sur $\operatorname{Div}(Y)$; la condition de (21.4.9) signifie que
$D'$ est de la forme $\operatorname{div}(t^n\circ f)$ pour un
$n\in\mathbb Z$. Cela montre aussitôt que cette condition entraîne la
condition 2\textdegree{} de l'énoncé. Inversement, supposons cette
dernière vérifiée (ainsi que la condition 1\textdegree{}); posant pour
simplifier $n=n_y$, tout revient à voir que l'on a alors
$D'=\operatorname{div}(t^n\circ f)$. Si $\eta$ est le point générique
de $Y$, les germes de $D'$ et de $\operatorname{div}(t^n\circ f)$ en
un point quelconque de $X_\eta$ sont tous deux nuls en vertu de
1\textdegree{}. De même, les germes de $D'$ et
$\operatorname{div}(t^n\circ f)$ en un point maximal de $X_y$ sont
égaux par la condition 2\textdegree{}. Mais en un point $x\in X_y$ qui
n'est pas maximal, on a
$\operatorname{prof}(\sh{O}_{X_y,x})\geq1$ puisque $X_y$ n'a pas de
cycle premier associé immergé, donc aussi
$\operatorname{prof}(\sh{O}_{X,x})\geq2$ par
{\hyperlink{ega4.c04.prop.6-3-1}{(6.3.1)}}. On a donc
$D'=\operatorname{div}(t^n\circ f)$ en vertu de
{\hyperlink{ega4.c04.prop.21-1-8}{(21.1.8)}}.
\end{proof}

\begin{corollary}[21.4.11]
\label{IV.21.4.11}
Supposons vérifiées les hypothèses générales de
{\hyperlink{ega4.c04.prop.21-4-9}{(21.4.9)}}, et supposons en outre que
$Y$ soit normal et que pour tout $y\in Y$ tel que
$\dim(\sh{O}_{Y,y})=1$, la fibre $X_y$ soit intègre. Alors, pour qu'un
diviseur $D'$ soit de la forme $f^*(D)$, où
$D\in\operatorname{Div}(Y)$, il faut et il suffit que, pour tout point
maximal $y$ de $Y$, l'image réciproque de $D'$ dans
$\operatorname{Div}(X_y)$ soit nulle.
\end{corollary}

\begin{proof}
Il suffit de vérifier que la condition 2\textdegree{} de
{\hyperlink{ega4.c04.cor.21-4-10}{(21.4.10)}} est alors
automatiquement vérifiée. Or, si $\dim(\sh{O}_{Y,y})=1$ et si $X_y$ est
intègre, on a, en l'unique point maximal $x$ de $X_y$,
$\dim(\sh{O}_{X,x})=1$ en vertu de
{\hyperlink{ega4.c04.cor.6-1-3}{(6.1.3)}}; en outre, comme
$\sh{O}_{Y,y}$ est un anneau régulier et $\sh{O}_{X_y,x}$ un corps,
$\sh{O}_{X,x}$ est régulier
{\hyperlink{ega4.c04.prop.6-5-1}{(6.5.1)}}, donc un anneau de valuation
discrète, et si $t$ est une uniformisante de $\sh{O}_{Y,y}$, son image
$t_y$ dans $\sh{O}_{X,x}$ est une uniformisante de $\sh{O}_{X,x}$
($0$, 17.3.3); l'image réciproque de $D'$ dans
$\operatorname{Div}(\operatorname{Spec}(\sh{O}_{X,x}))$ est donc bien
de la forme

\oldpage[IV]{361}
$\operatorname{div}(t_y^n)$, ce qui n'est autre que la condition
2\textdegree{} de {\hyperlink{ega4.c04.cor.21-4-10}{(21.4.10)}},
puisque $X_y$ n'a qu'un seul point maximal.
\end{proof}

\begin{remark}[21.4.12]
\label{IV.21.4.12}
Si dans {\hyperlink{ega4.c04.cor.21-4-11}{(21.4.11)}} on suppose
seulement que les $X_y$ sont réduits, il faut ajouter à la condition de
l'énoncé que les images réciproques de $D'$ dans chacun des ensembles
$\operatorname{Div}(\operatorname{Spec}(\sh{O}_{X,x}))$, où $x$
parcourt l'ensemble des points maximaux de $X_y$, sont toutes de la
forme $\operatorname{div}((t'_x)^n)$ où $t'_x$ est une uniformisante
de $\sh{O}_{X,x}$, l'entier $n$ étant le même pour tous ces points.
\end{remark}

\begin{corollary}[21.4.13]
\label{IV.21.4.13}
Supposons vérifiées les conditions de
{\hyperlink{ega4.c04.cor.21-4-11}{(21.4.11)}}. Pour tout
$\sh{O}_X$-Module inversible $\sh{L}$ tel que, pour tout point maximal
$\eta$ de $Y$, l'image réciproque
\[
 \sh{L}_\eta=\sh{L}\otimes_{\sh{O}_X}\sh{O}_{X_\eta}
\]
soit isomorphe à $\sh{O}_{X_\eta}$, il existe un
$\sh{O}_Y$-Module inversible $\sh{M}$ et un isomorphisme de
$f^*(\sh{M})$ sur $\sh{L}$.
\end{corollary}

\begin{proof}
Il suffit de montrer qu'il existe une section méromorphe régulière $s$
de $\sh{L}$ au-dessus de $X$ telle que le diviseur $D'$ correspondant
à $(\sh{L},s)$ soit tel que $D'_x=0$ en tout point $x$ d'une fibre
$X_\eta$ d'un point maximal $\eta$ de $Y$. On appliquera alors
{\hyperlink{ega4.c04.cor.21-4-11}{(21.4.11)}} à $D'$, et la conclusion
résultera de {\hyperlink{ega4.c04.def.21-4-2}{(21.4.2)}}.

Notons que les points de $\operatorname{Ass}(\sh{O}_X)$ sont au-dessus
des points maximaux de $Y$ {\hyperlink{ega4.c04.prop.3-3-1}{(3.3.1)}};
comme $Y$ est réduit et localement noethérien, il y a un ouvert dense
$V$ de $Y$, réunion d'une famille disjointe d'ouverts affines contenant
chacun un seul point maximal de $Y$, et $f^{-1}(V)$ est
schématiquement dense dans $X$
{\hyperlink{ega4.c04.prop.5-10-2}{(5.10.2)}}; on peut donc se borner au
cas où $Y$ est affine et intègre. Puisque, pour l'unique point maximal
$\eta$ de $Y$, $X_\eta$ est alors intègre,
$\operatorname{Ass}(\sh{O}_X)$ est réduit à un seul point $z$, qui est
le point générique de $X$. Soit $s_\eta$ la section de $\sh{L}_\eta$
au-dessus de $X_\eta$ qui correspond par isomorphie à la section unité
de $\sh{O}_{X_\eta}$, et soient $U$ un voisinage ouvert affine d'un
point $x\in X_\eta$ dans $X$ et $s$ une section de $\sh{L}$ au-dessus
de $U$ ayant même germe que $s_\eta$ au point $x$. Comme $U$ est
schématiquement dense dans $X$
({\hyperlink{ega4.c04.statement.20-2-13}{20.2.13}}, (iv)), $s$
s'identifie à une section méromorphe régulière bien déterminée de
$\sh{L}$ au-dessus de $X$, et il est clair
({\hyperlink{ega4.c04.prop.20-2-11}{20.2.11}} et
{\hyperlink{ega4.c04.prop.20-3-5}{20.3.5}}) que $s$ induit sur $X_\eta$
une section méromorphe régulière de $\sh{L}_\eta$ égale à $s_\eta$, et
répond donc à la question.
\end{proof}

\bigskip
\begin{center}
\small
1967. --- Imprimerie des Presses Universitaires de France. --- Vendôme
(France)\\
ÉDIT. N\textsuperscript{o} 29 470\qquad
IMPRIMÉ EN FRANCE\qquad
IMP. N\textsuperscript{o} 20 298
\end{center}
